% =============================================================================
% bab4-hasil.tex — BAB IV HASIL PENELITIAN DAN PEMBAHASAN
% Skripsi / Tesis / Disertasi (Lampiran I)
% =============================================================================

\chapter{HASIL PENELITIAN DAN PEMBAHASAN}
\label{bab:hasil}

\section{Hasil Penelitian}
\label{sec:hasil-penelitian}

\subsection{Gambaran Umum Objek Penelitian}
\label{subsec:gambaran-umum}

Uraikan gambaran umum objek/subjek penelitian di sini. Jelaskan profil singkat
perusahaan/organisasi/responden yang menjadi objek penelitian.

\subsection{Karakteristik Responden}
\label{subsec:karakteristik-responden}

Berikut adalah karakteristik responden yang menjadi sampel dalam penelitian ini:

\begin{table}[H]
\caption{Karakteristik Responden Berdasarkan Jenis Kelamin}
\label{tab:karakteristik-jenis-kelamin}
\begin{isitable}
\begin{tabular}{p{0.5cm} p{4cm} p{3cm} p{3cm}}
\toprule
\textbf{No} & \textbf{Jenis Kelamin} & \textbf{Jumlah} & \textbf{Persentase} \\
\midrule
1 & Laki-laki & 75 & 62,5\% \\
2 & Perempuan & 45 & 37,5\% \\
\midrule
& \textbf{Total} & \textbf{120} & \textbf{100\%} \\
\bottomrule
\end{tabular}
\end{isitable}
\end{table}

\subsection{Hasil Uji Validitas dan Reliabilitas}
\label{subsec:uji-valid-reliabel}

Sebelum melakukan analisis utama, terlebih dahulu dilakukan uji validitas dan
reliabilitas instrumen penelitian. Hasil uji validitas menunjukkan bahwa seluruh
item pernyataan memiliki nilai $r_{hitung} > r_{tabel}$ (0,179), sehingga semua
item dinyatakan valid.

Hasil uji reliabilitas menunjukkan koefisien Cronbach's Alpha sebesar \ldots,
yang berarti instrumen penelitian ini reliabel karena $\alpha \geq 0,60$.

\subsection{Analisis Deskriptif}
\label{subsec:analisis-deskriptif}

Berikut adalah hasil analisis deskriptif variabel-variabel penelitian:

\begin{table}[H]
\caption{Hasil Analisis Deskriptif}
\label{tab:analisis-deskriptif}
\begin{isitable}
\begin{tabular}{p{0.5cm} p{3.5cm} p{2cm} p{2cm} p{2cm} p{2cm} p{2cm}}
\toprule
\textbf{No} & \textbf{Variabel} & \textbf{N} & \textbf{Min} & \textbf{Max} & \textbf{Mean} & \textbf{Std. Dev} \\
\midrule
1 & Variabel X & 120 & 1,00 & 5,00 & 3,75 & 0,85 \\
2 & Variabel Y & 120 & 1,00 & 5,00 & 3,62 & 0,91 \\
\bottomrule
\end{tabular}
\end{isitable}
\end{table}

\subsection{Hasil Analisis [Regresi/SEM/dsb.]}
\label{subsec:hasil-analisis}

Uraikan hasil pengujian hipotesis dan analisis utama di sini. Sajikan hasil dalam
bentuk tabel dan/atau gambar yang mudah dipahami.

\begin{table}[H]
\caption{Hasil Pengujian Hipotesis}
\label{tab:hasil-hipotesis}
\begin{isitable}
\begin{tabular}{p{0.5cm} p{3cm} p{2cm} p{2cm} p{2cm} p{3cm}}
\toprule
\textbf{H} & \textbf{Hubungan} & \textbf{Koef.} & \textbf{t-stat} & \textbf{p-value} & \textbf{Kesimpulan} \\
\midrule
H\textsubscript{1} & X → Y & 0,45 & 5,23 & 0,000 & Diterima \\
H\textsubscript{2} & Z → Y & 0,32 & 3,87 & 0,001 & Diterima \\
\bottomrule
\end{tabular}
\end{isitable}
\end{table}

\section{Pembahasan}
\label{sec:pembahasan}

\subsection{Pembahasan Hipotesis 1}
\label{subsec:pembahasan-h1}

Berdasarkan hasil pengujian, hipotesis pertama diterima/ditolak. Hal ini menunjukkan
bahwa \ldots. Temuan ini konsisten/bertentangan dengan penelitian \parencite{contohSumber2025}
yang menyatakan bahwa \ldots

\subsection{Pembahasan Hipotesis 2}
\label{subsec:pembahasan-h2}

Berdasarkan hasil pengujian, hipotesis kedua diterima/ditolak. Hasil ini mengindikasikan
bahwa \ldots. Temuan ini sejalan dengan teori \ldots yang dikemukakan oleh \ldots

\section{Implikasi Manajerial}
\label{sec:implikasi}

Berdasarkan hasil penelitian dan pembahasan, terdapat beberapa implikasi manajerial
yang dapat disampaikan:

\begin{enumerate}
  \item Implikasi pertama bagi praktisi/manajemen.
  \item Implikasi kedua bagi kebijakan.
  \item Implikasi ketiga bagi pengembangan ilmu.
\end{enumerate}
